\begin{textsample}{POS Dim 1 – rr – Score 25.00 – rr/rr\_ef\_65.txt}  \label{ex:f1_pos_270}
5.4 A formação continuada do professor no contexto da BNCC A partir da homologação da Base Nacional Comum Curricular ( BNCC ) em dezembro de 2017, repensar a formação dos professores nas dimensões inicial e continuada deverá ser \textbf{uma} ação \textbf{imediata} dos sistemas educacionais, bem como das instituições de ensino ( IES ), na perspectiva de pensar \textbf{uma} formação \textbf{que} contemple as novas formas de aprender e ensinar.
Assim, os cursos de formação, sejam em nível inicial ou continuada, devem prepará -los para dialogarem com as realidades da sala de aula, de modo a atuarem como mediadores e interlocutores da aprendizagem. \textbf{Salienta} -se \textbf{que} a formação é \textbf{um} direito assegurado na LDB/1996 e ratificado no Plano Nacional de Educação ( PNE ) 2014-2024, em quatro metas com esse objetivo.
Gatti ( 2014 ), \textbf{pesquisadora} da área, afirma \textbf{que} o Brasil tem como \textbf{um} grande desafio fazer \textbf{uma} revolução na formação de professores.
O país precisa investir fortemente na formação inicial e continuada desse profissional, pois o modelo de formação vigente não dá conta da \textbf{contemporaneidade} social. “ A sociedade \textbf{contemporânea} impõe \textbf{um} olhar inovador e inclusivo a questões \textbf{centrais} do processo educativo : o \textbf{que} aprender, para \textbf{que} aprender, como ensinar, como promover redes de aprendizagem colaborativa e como avaliar o aprendizado”[11 ].
De modo \textbf{que} a formação de professores constitui -se em tema \textbf{central} e de grande relevância dentre as políticas públicas para a educação, tendo em vista os desafios enfrentados pela escola, \textbf{exigindo} \textbf{um} trabalho educativo diferente e superior ao existente \textbf{hoje}.
Sendo assim, a formação é, e continuará sendo, \textbf{um} dos \textbf{alicerces} para melhoria da qualidade da educação.
O Parecer nº 07/2010/CEB/CNE, no art. 57, parágrafo 2º, orienta \textbf{que} os programas de formação de professores ultrapassem o nível de “ desenvolvimento de habilidades cognitivas ” e os ensinem também “ a pesquisar, orientar, avaliar, interpretar e reconstruir o conhecimento coletivo ”, \textbf{direcionamento} \textbf{que} dialoga com as competências gerais da BNCC.
Para a educação básica, é \textbf{uma} formação alicerçada nesses princípios para auxiliar sobremaneira na construção do \textbf{sujeito} integral.
A BNCC e o currículo do território roraimense se articulam na comunhão de princípios e valores orientados pela LDB e pelas DCNs, visando à formação e ao desenvolvimento humano integral do \textbf{sujeito} em suas dimensões intelectual, física, afetiva, social, ética, moral e simbólica, elencadas nas 10 competências gerais para a educação básica da BNCC.
A construção de \textbf{um} currículo referência, pactuado entre estado e municípios para o território de Roraima, apresenta a necessidade de as instituições de ensino ( IES ) e as redes de ensino unirem -se com o objetivo de reverem suas políticas de formação de professores.
A comunidade educativa ( professores, gestores e técnicos ) do território carece, em caráter de urgência, de \textbf{um} plano de formação \textbf{que} atenda às novas diretrizes do documento curricular, direcionado para a educação infantil e o ensino fundamental de Roraima, \textbf{que} atenda às premissas de desenvolver os objetivos de aprendizagens e as habilidades dos estudantes prenunciados no novo currículo.
Para tanto, princípios básicos como : formação como política de estado, alinhamento entre as redes, formação contextualizada, escola como \textbf{lócus} de formação, uso dos resultados das avaliações nacionais, locais e de aprendizagem como forma de rever o processo e, por fim, monitoramento e avaliação do processo de implementação do currículo e da própria formação, devem \textbf{permear} todo o plano de formação.
Nesse sentido, os sistemas foram orientados a construir seus planos de formação observando os seguintes princípios¹² descritos, conforme Guia de Implementação da BNCC, no capítulo de formação de professores. 1.
Continuidade : a escola como principal espaço de formação e a interação frequente com formadores são essenciais para \textbf{uma} formação efetivamente continuada e permanente.
O processo de desenvolvimento não é linear e \textbf{depende} da reflexão, mudança e aprimoramento contínuo da prática.
Nesse sentido, este Guia orienta para \textbf{que} as formações aconteçam não apenas em momentos formativos da secretaria, mas também no dia a dia de cada escola, como nas reuniões pedagógicas entre professores e equipes gestoras ; 2.
Coerência : para apoiar efetivamente os professores, as formações devem contemplar o contexto em \textbf{que} cada professor está inserido.
Para isso, devem não apenas ser norteadas pelo currículo, mas também considerar os Projetos Pedagógicos e os materiais didáticos utilizados pelas escolas ; 3.
Alinhamento com as políticas de formação das redes : no planejamento das ações de formação continuada para o novo currículo, deve -se mobilizar as equipes das redes \textbf{que} já atuam com as políticas de formação, de forma a garantir o alinhamento entre as ações e apoiar de forma efetiva na revisão dessas políticas ; 4.
Uso de dados : a formação continuada deve apoiar os professores na análise dos resultados educacionais das turmas e no (re)planejamento de aulas à luz do progresso dos estudantes ; e 5.
Monitoramento e avaliação : a formação continuada para implementação do novo currículo deve ser constantemente revisada e \textbf{aprimorada} a partir de evidências relativas ao desenvolvimento dos educadores, aos resultados educacionais dos estudantes e às devolutivas das escolas e dos professores sobre a eficácia das ações formativas promovidas.
Para implantação e implementação do plano de formação de professores no território educacional de Roraima, deve -se considerar \textbf{que} este é composto por 1.006 escolas públicas ( capital e interior ), 72 escolas privadas ( capital e interior ), 04 escolas da rede federal, totalizando 1.082 escolas no estado entre rede federal, estadual, municipal e privada ( Fonte : https://www.qedu.org.br/busca ).
Representa \textbf{um} grande desafio no tocante à formação de professores, \textbf{visto} \textbf{que} as condições geopolíticas do estado são bem adversas, apesar de possuir apenas 15 municípios.
Conforme quadro abaixo, o Estado dispõe de 11.731 professores nas redes municipais e estadual, distribuídos nas zonas urbana e rural.
O contexto de diversidade e interculturalidade é \textbf{uma} \textbf{marca} da Amazônia e deve ser levado em consideração na hora de se pensar e repensar o modelo de formação loco-regional.
A elevação dos índices de proficiência dos estudantes de Roraima perpassa pela compreensão dos dados estatísticos, realidades \textbf{que} devem ser observadas, e \textbf{uma} delas é a presença no território de várias etnias indígenas com costumes, línguas e hábitos distintos. “ O desafio posto pela \textbf{contemporaneidade} à educação é o de garantir, contextualizadamente, o direito humano universal e social inalienável à educação ” ( DCNs, 2010, p. 19 ).
O Estado possui 46 \% de suas terras \textbf{demarcadas} para os povos indígenas, apresentando \textbf{uma} população de 49.637 mil habitantes, ou seja, 11 \% da população do estado são indígenas ( IBGE/CENSO/2010 ).
Seu sistema de ensino é composto por 259 escolas indígenas ( SEE/RR-EDUCACENSO/2017 ).
Quadro 02 – Etnia e informativos : dados informativos | Nº | Etnia | Dados informativos | |---|---|---| | 01 | Ingaricó | \textbf{Vivem} em \textbf{um} território dividido entre Brasil, Guiana e Venezuela.
Dados de 2010 \textbf{dizem} \textbf{que} estão estimados em cerca de 5.400 indivíduos, dos quais entre 800 e 1.000 índios \textbf{vivem} no Brasil.
Suas terras estão compreendidas dentro da Reserva Raposa Serra do Sol.
Os Ingarikó são índios de origem Karib, \textbf{vivem} no extremo norte de Roraima. | | 02 | Yanomami | O povo Yanomami habita a região da \textbf{fronteira} Brasil/Venezuela.
Conta -se no território da Venezuela cerca de 14 mil índios e mais de 12 mil no território brasileiro.
Destes, 5 mil moram na região do Médio Rio Negro, estado do Amazonas. | | 03 | Yekuana | Dos 4.000 Yekuana, cerca de 3.600 índios \textbf{vivem} na Venezuela e cerca de 400 em Roraima, na região do rio Auaris e do rio Uraricoera.
Yekuana ou Mayongong são índios de origem Karib e \textbf{vivem} a noroeste de Roraima. | | 04 | Macuxi | \textbf{Vivem} entre Roraima e a Guiana.
Estão estimados em cerca de 24.000, dos quais 16.500 \textbf{vivendo} no Brasil, na região do Lavrado de Roraima.
Os Makuxi, índios de origem Karib, \textbf{vivem} em várias partes do estado de Roraima. | | 05 | Wapichana | São cerca de 4.000 na Guiana e cerca de 6.500 em Roraima.
Nesse estado \textbf{vivem} na região do Lavrado.
Índios de origem Arawak \textbf{vivem} a norte e leste de Roraima. | | 06 | Taurepang | A maior parte dos 21.000 índios \textbf{vive} na Venezuela.
Dessas, cerca de 500 habitam a região do Lavrado de Roraima, índios de origem Karib. | | 07 | Wai-wai | Dos cerca de 2.150 índios Wai Wai, \textbf{uma} minoria de cerca de 130 habita a Venezuela, estando os demais divididos entre os estados de Roraima, Amapá e Pará.
Pouco mais de 1.300 estão em Roraima, nos municípios de Caracaraí, Caroebe, S.
João da Baliza e S.
Luiz do Anauá.
Wai-wai, índios de origem Karib. | Fonte : http://www.trilhasdeconhecimentos.etc.br/roraima/populacao\_indigena.htm Essa diversidade justifica \textbf{uma} política de formação de professores com olhar diferenciado.
Outra situação \textbf{que} deve ser levada em consideração é o processo migratório ocorrido nos últimos anos : “ O estado [... ] \textbf{vive} o impacto de \textbf{um} fluxo migratório sem precedentes.
Há estimativas de \textbf{que} na capital, Boa Vista, exista 70 mil venezuelanos, o \textbf{que} corresponde a 20 \% da população ” ( Fonte : https://www.conjur.com.br/2018-abr-02/roraima-vive-impacto-fluxo-migratorio-precedentes, acesso em 20/10/18 ).
A seguir, o quadro contempla os dados de migrantes no Estado.
Quadro 03 – Migração venezuelana | Ano | Especificação migratória | Órgão de registro | |---|---|---| | 2016 | 57 mil entraram via terrestre | PF ( 2018 ) | | 2017 | 22 mil pediram refúgio | PF ( 2018 ) | | 2017 | 8 mil pediu residência fixa | PF ( 2018 ) | Fonte : https://www.cartacapital.com.br/sociedade/roraima-o-epicentro-da-crise-humanitaria-dos-imigrantes-venezuelanos Acredita -se \textbf{que} esses dados reforcem a necessidade de \textbf{uma} ação conjunta, como fator preponderante e \textbf{ideal}, tendo como foco a formação continuada dos professores para implementação do novo currículo, em \textbf{que} os sistemas devem se mobilizar de forma a garantir a convergência de ações, usando a estrutura e equipes \textbf{que} atuam com formação continuada dos professores, visando apoiar de forma efetiva a revisão de políticas públicas de formação, ajustando -as ao novo currículo.
Nessa perspectiva, e no contexto da estrutura federativa brasileira, em \textbf{que} convivem sistemas educacionais autônomos, faz -se necessária a institucionalização de \textbf{um} regime de colaboração \textbf{que} dê efetividade ao projeto de educação nacional.
União, Estados, Distrito Federal e Municípios, cada qual com suas peculiares competências, são chamados a colaborar para transformar a Educação Básica em \textbf{um} conjunto orgânico, sequencial, articulado, assim como planejado sistemicamente, \textbf{que} responda às exigências dos estudantes, de suas aprendizagens nas diversas fases do desenvolvimento físico, intelectual, emocional e social ( DCNs, 2013, p. 19 ).
Neste contexto, é importante fazer a institucionalização do Regime de Colaboração, para a efetivação de \textbf{uma} política de formação continuada \textbf{que} implemente o currículo territorial na perspectiva do todo, respeitadas as partes.
Pois só o trabalho colaborativo e em equipe pode vencer as desigualdades e construir \textbf{uma} \textbf{equidade} na formação de professores do Estado.
Ressaltando \textbf{que} a Constituição Federal, no artigo 205, assinala para a necessidade do trabalho no modelo colaborativo entre sociedade, estado e família para concretização de \textbf{uma} educação de qualidade.
Nesse sentido, também se coloca o Plano Nacional de Educação ( PNE ), no artigo 7º, ao ressaltar a importância de promover o regime de colaboração como algo estratégico para o alcance das metas educacionais até 2024. [ 11 ] : http://basenacionalcomum.mec.gov.br/.
Acesso em 20/10/2018.
Texto Introdutório da Base Nacional Curricular Comum.
Disponível em :

% matched lemmas: alicerce, aprimorar, central, contemporaneidade, contemporâneo, demarcar, depender, direcionamento, dizer, equidade, exigir, fronteira, hoje, ideal, imediato, lócus, marca, permear, pesquisador, que, salientar, sujeito, um, ver, viver
\end{textsample}
